\documentclass[12pt, a4paper]{article}

% ── Codificación y lengua ──────────────────────────────────────────────────
\usepackage[utf8]{inputenc}
\usepackage[T1]{fontenc}
\usepackage[spanish]{babel}

% ── Tipografía y espaciado ─────────────────────────────────────────────────
\usepackage{lmodern}
\usepackage[protrusion=true,expansion=false]{microtype}
\usepackage{parskip}
\setlength{\parskip}{8pt}

% ── Márgenes ──────────────────────────────────────────────────────────────
\usepackage[top=2.5cm, bottom=2.5cm, left=2.8cm, right=2.8cm]{geometry}

% ── Color ─────────────────────────────────────────────────────────────────
\usepackage[dvipsnames, table]{xcolor}

% Paleta corporativa
\definecolor{primario}{HTML}{1D9E75}     % teal principal
\definecolor{primarioClaro}{HTML}{E1F5EE}
\definecolor{secundario}{HTML}{534AB7}   % púrpura
\definecolor{secundarioClaro}{HTML}{EEEDFE}
\definecolor{acento}{HTML}{D85A30}       % coral
\definecolor{acentoClaro}{HTML}{FAECE7}
\definecolor{ambar}{HTML}{BA7517}
\definecolor{ambarClaro}{HTML}{FAEEDA}
\definecolor{gris}{HTML}{5F5E5A}
\definecolor{grisClaro}{HTML}{F1EFE8}
\definecolor{grisMedio}{HTML}{D3D1C7}

% ── Cabeceras y pies de página ─────────────────────────────────────────────
\usepackage{fancyhdr}
\pagestyle{fancy}
\fancyhf{}
\renewcommand{\headrulewidth}{0.4pt}
\fancyhead[L]{\small\textcolor{gris}{App CAA para TEA --- Blueprint Técnico}}
\fancyhead[R]{\small\textcolor{gris}{\thepage}}
\fancyfoot[C]{\small\textcolor{grisMedio}{Documento confidencial · 2025}}

% ── Secciones ─────────────────────────────────────────────────────────────
\usepackage{titlesec}
\titleformat{\section}
  {\large\bfseries\color{primario}}
  {\thesection.}{0.8em}{}[\vspace{-2pt}\textcolor{grisMedio}{\rule{\linewidth}{0.6pt}}]

\titleformat{\subsection}
  {\normalsize\bfseries\color{secundario}}
  {\thesubsection.}{0.7em}{}

\titleformat{\subsubsection}
  {\normalsize\bfseries\color{gris}}
  {\thesubsubsection.}{0.6em}{}

% ── Cajas de contenido ────────────────────────────────────────────────────
\usepackage[most]{tcolorbox}

% Caja de "punto clave"
\newtcolorbox{puntoclave}[1][]{
  colback=primarioClaro, colframe=primario,
  boxrule=0.5pt, arc=6pt,
  left=10pt, right=10pt, top=8pt, bottom=8pt,
  fonttitle=\small\bfseries\color{primario},
  title={#1}
}

% Caja técnica (código / esquema)
\newtcolorbox{cajatecnica}[1][]{
  colback=grisClaro, colframe=grisMedio,
  boxrule=0.4pt, arc=4pt,
  left=10pt, right=10pt, top=8pt, bottom=8pt,
  fonttitle=\small\bfseries\color{gris},
  title={#1}
}

% Caja de módulo de la app
\newtcolorbox{cajamodulo}[2][]{
  colback=#1, colframe=#1!70!black,
  boxrule=0.5pt, arc=8pt,
  left=12pt, right=12pt, top=10pt, bottom=10pt,
  title={\textbf{#2}},
  fonttitle=\normalsize\bfseries,
  #1
}

% Caja de alerta / nota
\newtcolorbox{nota}[1][Nota]{
  colback=ambarClaro, colframe=ambar,
  boxrule=0.5pt, arc=5pt,
  left=10pt, right=8pt, top=6pt, bottom=6pt,
  fonttitle=\small\bfseries\color{ambar},
  title={#1}
}

% ── Listas ────────────────────────────────────────────────────────────────
\usepackage{enumitem}
\setlist[itemize]{leftmargin=1.4em, itemsep=3pt, topsep=4pt}
\setlist[enumerate]{leftmargin=1.4em, itemsep=3pt, topsep=4pt}

% ── Tablas ────────────────────────────────────────────────────────────────
\usepackage{booktabs}
\usepackage{array}
\usepackage{longtable}
\newcolumntype{L}[1]{>{\raggedright\arraybackslash}p{#1}}
\newcolumntype{C}[1]{>{\centering\arraybackslash}p{#1}}

% ── Hyperlinks ────────────────────────────────────────────────────────────
\usepackage{hyperref}
\hypersetup{
  colorlinks=true,
  linkcolor=secundario,
  urlcolor=primario,
  citecolor=gris,
  pdftitle={App CAA para TEA - Blueprint Técnico},
  pdfauthor={Equipo de Desarrollo},
  pdfsubject={Comunicación Aumentativa y Alternativa}
}

% ── TikZ para diagramas ───────────────────────────────────────────────────
\usepackage{tikz}
\usetikzlibrary{shapes.rounded, arrows.meta, positioning, fit, backgrounds}

% ── Código fuente ─────────────────────────────────────────────────────────
\usepackage{listings}
\lstset{
  basicstyle=\small\ttfamily\color{gris},
  backgroundcolor=\color{grisClaro},
  frame=single, framerule=0.4pt, rulecolor=\color{grisMedio},
  breaklines=true, breakatwhitespace=true,
  xleftmargin=10pt, xrightmargin=10pt,
  numbers=left, numberstyle=\tiny\color{grisMedio},
  numbersep=8pt,
  keywordstyle=\bfseries\color{secundario},
  commentstyle=\itshape\color{gris!70},
  stringstyle=\color{primario},
  tabsize=2, showstringspaces=false
}

% ── Graficados inline sencillos ───────────────────────────────────────────
\usepackage{graphicx}

% ─────────────────────────────────────────────────────────────────────────
\begin{document}

% ══════════════════════════════════════════════════════════════════════════
%  PORTADA
% ══════════════════════════════════════════════════════════════════════════
\begin{titlepage}
\thispagestyle{empty}
\begin{tikzpicture}[remember picture, overlay]
  % Banda superior
  \fill[primario] (current page.north west) rectangle
    ([yshift=-3.5cm]current page.north east);
  % Banda inferior delgada
  \fill[primarioClaro] (current page.south west) rectangle
    ([yshift=1.2cm]current page.south east);
  % Acento lateral
  \fill[acento] ([xshift=0cm]current page.north west) rectangle
    ([xshift=0.8cm, yshift=-3.5cm]current page.north west);
\end{tikzpicture}

\vspace*{1.8cm}
{\color{white}\fontsize{11}{14}\selectfont\textbf{BLUEPRINT TÉCNICO Y DE DISEÑO}}

\vspace{0.4cm}
{\color{white}\fontsize{24}{30}\selectfont\textbf{Aplicación CAA para TEA}}

\vspace{0.3cm}
{\color{white!80!primario}\fontsize{14}{18}\selectfont
Comunicación Aumentativa y Alternativa\\
para niños con Trastorno del Espectro Autista\\
Nivel 2--3, rango de edad 3--6 años}

\vspace{2.5cm}

\begin{tcolorbox}[colback=white, colframe=grisMedio, arc=8pt,
  boxrule=0.5pt, left=16pt, right=16pt, top=12pt, bottom=12pt,
  width=\textwidth]
\begin{tabular}{L{4cm} L{9cm}}
\textcolor{gris}{\small\textbf{Tipo de documento}} &
  \small Blueprint técnico, de datos y de diseño \\[4pt]
\textcolor{gris}{\small\textbf{Audiencia}} &
  \small Evaluadores académicos · Inversores · Equipo técnico · Profesionales sanitarios/educativos \\[4pt]
\textcolor{gris}{\small\textbf{Versión}} &
  \small 1.0 --- Documento inicial \\[4pt]
\textcolor{gris}{\small\textbf{Fecha}} &
  \small 2025 \\[4pt]
\textcolor{gris}{\small\textbf{Estado}} &
  \small \textcolor{acento}{\textbf{Pre-desarrollo}} --- Fase de diseño y planificación \\
\end{tabular}
\end{tcolorbox}

\vfill

{\small\textcolor{gris}{Documento preparado para presentación y revisión. No contiene código funcional final.\\
Todos los esquemas, algoritmos y decisiones técnicas son previos a la implementación.}}

\end{titlepage}

% ══════════════════════════════════════════════════════════════════════════
%  ÍNDICE
% ══════════════════════════════════════════════════════════════════════════
\newpage
\tableofcontents
\newpage

% ══════════════════════════════════════════════════════════════════════════
%  1. RESUMEN EJECUTIVO
% ══════════════════════════════════════════════════════════════════════════
\section{Resumen ejecutivo}

Este documento describe el diseño completo de una aplicación de \textbf{Comunicación Aumentativa y Alternativa (CAA)} orientada a niños de entre 3 y 6 años con diagnóstico de \textbf{Trastorno del Espectro Autista (TEA) en niveles 2 y 3}.

La aplicación no busca sustituir el lenguaje oral, sino ofrecer un andamiaje estructurado para que el niño pueda:

\begin{itemize}
  \item \textbf{Construir frases} combinando pictogramas de forma secuencial.
  \item \textbf{Anticipar rutinas} mediante representaciones visuales claras del tiempo.
  \item \textbf{Regular sus emociones} identificando estados internos y estrategias de calma.
\end{itemize}

Paralelamente, la plataforma proporciona a padres y docentes un \textbf{panel de seguimiento} con indicadores objetivos del progreso comunicativo del niño.

\begin{puntoclave}[Propuesta de valor]
La aplicación combina intervención terapéutica basada en evidencia (CAA, pictogramas ARASAAC) con tecnología moderna (PWA, TTS nativo, API REST), generando datos clínicamente relevantes sin añadir carga a las familias.
\end{puntoclave}

\subsection{Población objetivo}

\begin{tabular}{L{4cm} L{10cm}}
\toprule
\textbf{Criterio} & \textbf{Descripción} \\
\midrule
Edad & 3 a 6 años \\
Diagnóstico & TEA nivel 2--3 (DSM-5 / CIE-11) \\
Perfil lingüístico & Verbales o con comunicación funcional emergente \\
Contexto de uso & Domicilio familiar, aula ordinaria o de apoyo, centros de atención temprana \\
Dispositivo esperado & Tablet (iOS o Android), resolución mínima 768×1024 px \\
\bottomrule
\end{tabular}

\subsection{Objetivos terapéuticos}

\begin{enumerate}
  \item Aumentar la \textbf{longitud media de frase} (LMF) en comunicación espontánea.
  \item Ampliar el \textbf{vocabulario funcional} (diversidad léxica de pictogramas usados).
  \item Fomentar el uso \textbf{autónomo} del regulador emocional sin mediación adulta.
  \item Reducir la ansiedad ante cambios de rutina mediante la \textbf{anticipación visual}.
\end{enumerate}

% ══════════════════════════════════════════════════════════════════════════
%  2. CONTEXTO Y JUSTIFICACIÓN
% ══════════════════════════════════════════════════════════════════════════
\section{Contexto y justificación}

\subsection{¿Qué es la Comunicación Aumentativa y Alternativa?}

La CAA engloba todos los métodos, herramientas y estrategias que complementan o reemplazan el habla cuando esta es insuficiente para satisfacer las necesidades comunicativas de una persona. Los sistemas de pictogramas --- representaciones gráficas simples asociadas a palabras --- constituyen la herramienta más extendida en intervención con TEA.

El portal \textbf{ARASAAC} (Centro Aragonés para la Comunicación Aumentativa y Alternativa), dependiente del Gobierno de Aragón, ofrece más de 40\,000 pictogramas bajo licencia Creative Commons, con API pública gratuita. Es el estándar de facto en contextos hispanohablantes.

\subsection{Evidencia sobre eficacia}

Numerosos estudios respaldan el uso de sistemas de CAA basados en pictogramas en niños con TEA:

\begin{itemize}
  \item Mejoran la iniciación comunicativa y reducen conductas disruptivas asociadas a la frustración por no poder expresarse.
  \item La anticipación visual de rutinas reduce significativamente la ansiedad ante transiciones.
  \item Los sistemas de regulación emocional basados en pictogramas favorecen la identificación y el etiquetado de emociones propias.
\end{itemize}

\begin{nota}[Sobre el nivel de TEA]
Los niveles 2 y 3 implican necesidades de apoyo sustancial. Estos niños pueden tener lenguaje oral, pero con dificultades marcadas en la comunicación social espontánea, rigidez ante cambios y alta variabilidad en sus estrategias de autorregulación. La app se diseña específicamente para este perfil, no para TEA nivel 1 ni para usuarios sin lenguaje funcional alguno.
\end{nota}

\subsection{Brecha de mercado}

Las soluciones existentes presentan limitaciones relevantes:

\begin{tabular}{L{4cm} L{5cm} L{5cm}}
\toprule
\textbf{Herramienta} & \textbf{Fortaleza} & \textbf{Limitación} \\
\midrule
Proloquo2Go & Potente y completa & Coste elevado, complejidad de configuración \\
LetMeTalk & Gratuita y open source & Sin panel de seguimiento, UX poco adaptada \\
Pictogramas ARASAAC & Recurso libre de calidad & No es una app, solo banco de imágenes \\
Tableros físicos & Bajo coste, sin pantalla & Sin registro, sin adaptación dinámica \\
\bottomrule
\end{tabular}

\vspace{6pt}
Nuestra propuesta ocupa el espacio de una herramienta \textbf{gratuita o de bajo coste, en español, con seguimiento clínico integrado} y diseño específico para el perfil TEA nivel 2--3.

% ══════════════════════════════════════════════════════════════════════════
%  3. ARQUITECTURA TÉCNICA
% ══════════════════════════════════════════════════════════════════════════
\section{Arquitectura técnica}

\subsection{Stack tecnológico}

\begin{tabular}{L{3.5cm} L{4cm} L{6.5cm}}
\toprule
\textbf{Capa} & \textbf{Tecnología} & \textbf{Justificación} \\
\midrule
Frontend & React + Vite (PWA) & Instalable en tablet sin app store, offline-first \\
Backend & Node.js + Express & Ecosistema amplio, bajo coste de hosting \\
Base de datos & PostgreSQL & Integridad relacional, agregaciones SQL para KPIs \\
Pictogramas & API ARASAAC & 40\,000+ pictogramas, licencia CC, gratuita \\
Síntesis de voz & Web Speech API & Nativa en navegador, sin coste ni latencia de red \\
Hosting & Railway / Render & Tier gratuito viable para MVP \\
\bottomrule
\end{tabular}

\subsection{Diagrama de arquitectura}

\begin{center}
\begin{tikzpicture}[
  node distance=0.8cm and 1.2cm,
  box/.style={rounded corners=6pt, minimum width=3cm, minimum height=1cm,
              text centered, font=\small\bfseries, draw, thick},
  arrow/.style={-{Stealth[scale=1.1]}, thick, color=gris},
  label/.style={font=\footnotesize\itshape, color=gris}
]

% Dispositivo usuario
\node[box, fill=primarioClaro, draw=primario, color=primario] (tablet)
  {Tablet / Navegador\\{\footnotesize React PWA}};

% Backend
\node[box, fill=secundarioClaro, draw=secundario, color=secundario,
      right=2.2cm of tablet] (api)
  {API REST\\{\footnotesize Node.js + Express}};

% Base de datos
\node[box, fill=ambarClaro, draw=ambar, color=ambar,
      right=2.2cm of api] (db)
  {PostgreSQL\\{\footnotesize Base de datos}};

% ARASAAC
\node[box, fill=acentoClaro, draw=acento, color=acento,
      above=1.2cm of api] (arasaac)
  {API ARASAAC\\{\footnotesize Pictogramas CC}};

% TTS
\node[box, fill=grisClaro, draw=grisMedio, color=gris,
      below=1.2cm of tablet] (tts)
  {Web Speech API\\{\footnotesize TTS nativo}};

% Flechas
\draw[arrow] (tablet) -- node[above, label]{HTTP/JSON} (api);
\draw[arrow] (api) -- node[above, label]{SQL} (db);
\draw[arrow] (api) -- node[right, label]{GET pictogramas} (arasaac);
\draw[arrow, dashed] (arasaac) -- node[left, label]{URL imágenes} (tablet);
\draw[arrow, <->] (tablet) -- node[right, label]{síntesis voz} (tts);

\end{tikzpicture}
\end{center}

\subsection{Justificación de PostgreSQL vs MongoDB}

Se eligió \textbf{PostgreSQL} (base de datos relacional) por las siguientes razones:

\begin{itemize}
  \item Las relaciones entre entidades son fuertemente estructuradas: un usuario genera frases, cada frase contiene pictogramas en un orden específico.
  \item Los KPIs del panel adulto requieren \textbf{agregaciones y JOINs eficientes} (promedio semanal, conteo de valores distintos, agrupación por fecha).
  \item El historial clínico y educativo exige \textbf{integridad referencial real}: no puede existir una frase sin un usuario válido.
  \item El campo \texttt{pictogram\_ids} en la tabla de frases usa el tipo nativo \texttt{jsonb} de PostgreSQL, que permite almacenar el orden sin sacrificar la eficiencia en consulta.
\end{itemize}

% ══════════════════════════════════════════════════════════════════════════
%  4. ESTRUCTURA DE DATOS
% ══════════════════════════════════════════════════════════════════════════
\section{Estructura de datos (esquema de base de datos)}

\subsection{Entidades y relaciones}

El esquema consta de \textbf{seis tablas principales}. Las relaciones clave son:

\begin{itemize}
  \item \texttt{USERS} $\rightarrow$ \texttt{GENERATED\_PHRASES} (un usuario genera muchas frases)
  \item \texttt{USERS} $\rightarrow$ \texttt{EMOTIONAL\_LOGS} (un usuario registra muchas emociones)
  \item \texttt{USERS} $\rightarrow$ \texttt{USAGE\_EVENTS} (telemetría completa por usuario)
  \item \texttt{CATEGORIES} $\rightarrow$ \texttt{PICTOGRAMS} (una categoría agrupa muchos pictogramas)
\end{itemize}

\subsection{Definición de tablas}

\subsubsection{USERS --- Gestión de usuarios y roles}

\begin{cajatecnica}[Tabla: users]
\begin{lstlisting}[language=SQL]
CREATE TABLE users (
  id            UUID PRIMARY KEY DEFAULT gen_random_uuid(),
  name          VARCHAR(100) NOT NULL,
  role          VARCHAR(20) NOT NULL
                  CHECK (role IN ('child','parent','teacher')),
  linked_adult_id UUID REFERENCES users(id),
  avatar_color  VARCHAR(7),        -- hex color para avatar
  created_at    TIMESTAMPTZ NOT NULL DEFAULT NOW()
);
\end{lstlisting}
\end{cajatecnica}

El campo \texttt{role} define el tipo de usuario. Un niño puede tener un adulto vinculado principal (\texttt{linked\_adult\_id}); para relaciones muchos-a-muchos se añadiría una tabla intermedia \texttt{user\_links}.

\subsubsection{CATEGORIES y PICTOGRAMS --- Biblioteca de pictogramas}

\begin{cajatecnica}[Tablas: categories y pictograms]
\begin{lstlisting}[language=SQL]
CREATE TABLE categories (
  id            UUID PRIMARY KEY DEFAULT gen_random_uuid(),
  name          VARCHAR(60) NOT NULL,  -- 'Rutinas', 'Emociones'...
  color_code    VARCHAR(7),            -- color de la categoria en UI
  display_order INT NOT NULL DEFAULT 0
);

CREATE TABLE pictograms (
  id            UUID PRIMARY KEY DEFAULT gen_random_uuid(),
  arasaac_id    INT NOT NULL UNIQUE,   -- ID en la API de ARASAAC
  word          VARCHAR(80) NOT NULL,  -- palabra escrita asociada
  image_url     TEXT NOT NULL,         -- URL cacheada de ARASAAC
  category_id   UUID NOT NULL REFERENCES categories(id)
);
\end{lstlisting}
\end{cajatecnica}

\subsubsection{GENERATED\_PHRASES --- Historial de frases}

\begin{cajatecnica}[Tabla: generated\_phrases]
\begin{lstlisting}[language=SQL]
CREATE TABLE generated_phrases (
  id              UUID PRIMARY KEY DEFAULT gen_random_uuid(),
  user_id         UUID NOT NULL REFERENCES users(id),
  pictogram_ids   JSONB NOT NULL,
  -- Formato: [{"id": "uuid-picto", "position": 1}, ...]
  phrase_length   INT NOT NULL,        -- numero de pictogramas
  created_at      TIMESTAMPTZ NOT NULL DEFAULT NOW()
);

CREATE INDEX idx_phrases_user_date
  ON generated_phrases(user_id, created_at DESC);
\end{lstlisting}
\end{cajatecnica}

\subsubsection{EMOTIONAL\_LOGS --- Registro emocional}

\begin{cajatecnica}[Tabla: emotional\_logs]
\begin{lstlisting}[language=SQL]
CREATE TABLE emotional_logs (
  id                UUID PRIMARY KEY DEFAULT gen_random_uuid(),
  user_id           UUID NOT NULL REFERENCES users(id),
  emotion           VARCHAR(40) NOT NULL,
  intensity         SMALLINT CHECK (intensity BETWEEN 1 AND 3),
  strategy_chosen   VARCHAR(80),  -- NULL si no eligio estrategia
  created_at        TIMESTAMPTZ NOT NULL DEFAULT NOW()
);
\end{lstlisting}
\end{cajatecnica}

\subsubsection{USAGE\_EVENTS --- Telemetría y KPIs}

\begin{cajatecnica}[Tabla: usage\_events]
\begin{lstlisting}[language=SQL]
CREATE TABLE usage_events (
  id          UUID PRIMARY KEY DEFAULT gen_random_uuid(),
  user_id     UUID NOT NULL REFERENCES users(id),
  event_type  VARCHAR(60) NOT NULL,
  -- Ejemplos: 'phrase_built', 'emotion_logged',
  --           'anticipation_viewed', 'session_start'
  details     JSONB,
  created_at  TIMESTAMPTZ NOT NULL DEFAULT NOW()
);

CREATE INDEX idx_events_user_type
  ON usage_events(user_id, event_type, created_at DESC);
\end{lstlisting}
\end{cajatecnica}

% ══════════════════════════════════════════════════════════════════════════
%  5. SISTEMA DE SEGUIMIENTO Y KPIs
% ══════════════════════════════════════════════════════════════════════════
\section{Sistema de seguimiento: indicadores de progreso (KPIs)}

El panel adulto calcula cuatro indicadores a partir de los datos almacenados. A continuación se describe su lógica en pseudocódigo y la consulta SQL equivalente.

\subsection{KPI 1 --- Longitud media de frase (LMF)}

\begin{puntoclave}[¿Qué mide?]
El promedio de pictogramas por frase en una semana dada. Un aumento sostenido indica mayor complejidad comunicativa.
\end{puntoclave}

\begin{cajatecnica}[Pseudocódigo --- LMF semanal]
\begin{lstlisting}
FUNCTION calcLMF(userId, weekStart):
  phrases = DB.query(
    "SELECT phrase_length
     FROM generated_phrases
     WHERE user_id = $1
       AND created_at >= $2
       AND created_at <  $2 + INTERVAL '7 days'",
    userId, weekStart
  )
  IF phrases.length == 0: RETURN null
  RETURN ROUND(SUM(phrase.phrase_length) / phrases.length, 2)

-- Para la grafica historica (ultimas 8 semanas):
FUNCTION lmfSeries(userId, weeks = 8):
  RETURN [
    { week: w, lmf: calcLMF(userId, w) }
    FOR w IN last_n_week_starts(weeks)
  ]
\end{lstlisting}
\end{cajatecnica}

\subsection{KPI 2 --- Diversidad léxica}

\begin{puntoclave}[¿Qué mide?]
El número de pictogramas \emph{distintos} utilizados en los últimos 30 días. Refleja la amplitud del vocabulario funcional activo.
\end{puntoclave}

\begin{cajatecnica}[Consulta SQL --- Diversidad léxica]
\begin{lstlisting}[language=SQL]
SELECT COUNT(DISTINCT elem->>'id') AS unique_pictograms
FROM generated_phrases,
     jsonb_array_elements(pictogram_ids) AS elem
WHERE user_id    = $1
  AND created_at >= NOW() - INTERVAL '30 days';
\end{lstlisting}
\end{cajatecnica}

\subsection{KPI 3 --- Frecuencia de uso}

\begin{puntoclave}[¿Qué mide?]
El número de días activos por semana. Un uso frecuente (3--5 días/semana) es indicador de integración real de la herramienta en la vida diaria.
\end{puntoclave}

\begin{cajatecnica}[Consulta SQL --- Días activos por semana]
\begin{lstlisting}[language=SQL]
SELECT
  DATE_TRUNC('week', created_at)     AS week,
  COUNT(DISTINCT DATE(created_at))   AS active_days
FROM usage_events
WHERE user_id    = $1
  AND created_at >= NOW() - INTERVAL '8 weeks'
GROUP BY week
ORDER BY week;
\end{lstlisting}
\end{cajatecnica}

\subsection{KPI 4 --- Uso autónomo del regulador emocional}

\begin{puntoclave}[¿Qué mide?]
El porcentaje de registros emocionales en los que el niño seleccionó una estrategia de regulación (sin que el adulto la eligiese por él). Es el indicador más relevante de autorregulación.
\end{puntoclave}

\begin{cajatecnica}[Consulta SQL --- Tasa de autonomía emocional]
\begin{lstlisting}[language=SQL]
SELECT
  COUNT(*)                                  AS total_logs,
  COUNT(*) FILTER (WHERE strategy_chosen
                         IS NOT NULL)       AS autonomous_logs,
  ROUND(
    100.0 * COUNT(*) FILTER (
      WHERE strategy_chosen IS NOT NULL)
    / NULLIF(COUNT(*), 0), 1
  )                                         AS autonomy_rate
FROM emotional_logs
WHERE user_id    = $1
  AND created_at >= NOW() - INTERVAL '30 days';
\end{lstlisting}
\end{cajatecnica}

\subsection{Endpoint del dashboard}

Un único endpoint agrega los cuatro KPIs y los devuelve en una sola llamada, evitando múltiples peticiones desde el frontend:

\begin{cajatecnica}[]
\begin{lstlisting}
GET /api/dashboard/:childId?weeks=8&days=30

Respuesta JSON:
{
  "lmfSeries":         [{ week, lmf }, ...],
  "diversity":         { count, periodDays },
  "frequency":         [{ week, activeDays }, ...],
  "emotionalAutonomy": { total, autonomous, rate }
}
\end{lstlisting}
\end{cajatecnica}

% ══════════════════════════════════════════════════════════════════════════
%  6. DISEÑO DE INTERFAZ
% ══════════════════════════════════════════════════════════════════════════
\section{Diseño de interfaz: principios y módulos}

\subsection{Principios de diseño para TEA nivel 2--3}

El diseño de la interfaz parte de la investigación en necesidades perceptivas y cognitivas de niños con TEA. Los principios no negociables son:

\begin{tabular}{L{4cm} L{10cm}}
\toprule
\textbf{Principio} & \textbf{Aplicación concreta} \\
\midrule
Reducción de sobrecarga sensorial & Fondo neutro (blanco roto / gris muy claro), paleta pastel no saturada \\
Máximo 6 ítems en pantalla & Nunca más de 6 pictogramas visibles simultáneamente \\
Previsibilidad & Navegación lineal, sin menús ocultos ni gestos complejos \\
Asociación constante & Cada pictograma lleva su palabra escrita debajo, siempre \\
Pictogramas ARASAAC & Estándar reconocido, consistencia visual, licencia libre \\
Tamaño homogéneo & Todos los pictogramas del mismo tamaño, mínimo 80×80 px \\
Feedback inmediato & Cada interacción tiene respuesta visual y/o sonora clara \\
\bottomrule
\end{tabular}

\subsection{Paleta de colores}

\begin{tabular}{L{3cm} L{3cm} L{8cm}}
\toprule
\textbf{Módulo} & \textbf{Color} & \textbf{Uso} \\
\midrule
Frases & Verde-teal & Barra de frase, botón reproducir \\
Emociones & Coral/rosa suave & Grid de emociones, estrategias \\
Anticipación & Ámbar claro & Bloques de tiempo (Ahora / Después / Luego) \\
Panel adulto & Púrpura suave & Cabeceras, gráficas de progreso \\
Neutro & Gris muy claro & Fondo general, separadores \\
\bottomrule
\end{tabular}

\subsection{Módulo A: Construcción de frases}

\textbf{Descripción:} Permite al niño construir una frase combinando pictogramas de forma secuencial y reproducirla por voz.

\textbf{Flujo de interacción:}

\begin{enumerate}
  \item La pantalla se abre con la \textbf{barra de frase vacía} en la parte superior y el \textbf{grid de categorías} (máx. 6) debajo.
  \item El niño toca una \textbf{categoría} $\to$ el grid se actualiza mostrando los pictogramas de esa categoría (máx. 6).
  \item El niño toca un \textbf{pictograma} $\to$ aparece en la barra de frase con animación suave.
  \item Puede tocar un pictograma en la barra para \textbf{eliminarlo} (corrección sin pantallas nuevas).
  \item Al tocar el \textbf{botón ``Reproducir''} (icono grande, verde), el sistema sintetiza la frase por voz y la guarda en el historial.
  \item Un botón de \textbf{papelera} limpia la barra para empezar de nuevo.
\end{enumerate}

\textbf{Componentes React principales:}

\begin{cajatecnica}[]
\begin{lstlisting}
<PhraseBuilder>
  <PhraseBar />          {/* frases acumuladas, clicables para borrar */}
  <CategoryStrip />      {/* selector de categoria, maximo 6 pills */}
  <PictoGrid>
    <PictoCard />        {/* pictograma + palabra, maximo 6 por grid */}
  </PictoGrid>
  <PlayButton />         {/* TTS: Web Speech API, lang="es-ES" */}
  <ClearButton />
</PhraseBuilder>
\end{lstlisting}
\end{cajatecnica}

\textbf{Estado (useReducer):}
\texttt{\{ phrase: Picto[], activeCategory: string|null, isPlaying: boolean \}}

\subsection{Módulo B: Anticipación visual}

\textbf{Descripción:} Muestra al niño qué va a pasar en los próximos pasos de su jornada, reduciendo la ansiedad ante transiciones.

\textbf{Estructura visual:} Pantalla dividida en \textbf{tres bloques} de igual tamaño con fondos pastel diferenciados:

\begin{tabular}{C{3cm} C{3cm} C{3cm}}
\toprule
\textbf{AHORA} & \textbf{DESPUÉS} & \textbf{LUEGO} \\
\midrule
Pictograma grande & Pictograma grande & Pictograma grande \\
Fondo ámbar & Fondo teal claro & Fondo gris claro \\
\bottomrule
\end{tabular}

\vspace{8pt}
\textbf{Flujo de interacción:}

\begin{enumerate}
  \item El adulto configura los tres slots desde el panel adulto (con PIN de acceso).
  \item El niño visualiza la secuencia. Al completar ``Ahora'', el adulto toca el bloque $\to$ aparece un \textbf{check grande} y los bloques avanzan.
  \item Los datos se obtienen de \texttt{GET /api/schedule/:childId/today}.
\end{enumerate}

\subsection{Módulo C: Regulador emocional}

\textbf{Descripción:} Ayuda al niño a identificar su emoción y acceder a estrategias de regulación adecuadas.

\textbf{Flujo de tres pasos:}

\begin{enumerate}
  \item \textbf{Grid de emociones} (máx. 6): Feliz, Triste, Enfadado, Asustado, Sorprendido, Calmado. El niño toca la suya.
  \item \textbf{Intensidad} (opcional, configurable): tres niveles visuales mediante tamaño de círculo creciente.
  \item \textbf{Estrategias} (máx. 3 grandes): pictogramas de acciones calmantes según la emoción elegida (ej: Respirar, Zona tranquila, Música).
\end{enumerate}

Al seleccionar una estrategia, el sistema guarda \texttt{POST /api/emotional-log} con \texttt{\{emotion, intensity, strategy\_chosen\}} y muestra confirmación positiva. Si el niño sale sin elegir estrategia, el log se guarda con \texttt{strategy\_chosen = null} (uso no autónomo).

\subsection{Módulo D: Panel de seguimiento (vista adulto)}

\textbf{Descripción:} Interfaz exclusiva para padres y docentes, accesible con PIN, que muestra el progreso del niño de forma visual y comprensible.

\textbf{Estructura:}

\begin{enumerate}
  \item \textbf{Selector de niño}: dropdown si el adulto tiene varios perfiles vinculados.
  \item \textbf{4 KpiCards} en la parte superior: LMF actual, Diversidad léxica, Días activos, Tasa de autonomía emocional. Cada card incluye flecha de tendencia respecto a la semana anterior.
  \item \textbf{Gráficas} (usando Recharts):
    \begin{itemize}
      \item \emph{Línea suave}: LMF por semana (eje X = semanas, eje Y = pictogramas).
      \item \emph{Barras horizontales}: uso por categorías de pictogramas.
      \item \emph{Donut}: distribución de emociones registradas.
    \end{itemize}
  \item \textbf{Selector de rango}: 7 días / 30 días / 90 días. Reactivo: todos los charts se re-cargan.
\end{enumerate}

% ══════════════════════════════════════════════════════════════════════════
%  7. ÁRBOL DE COMPONENTES REACT
% ══════════════════════════════════════════════════════════════════════════
\section{Árbol de componentes React}

\begin{cajatecnica}[Estructura de componentes (simplificada)]
\begin{lstlisting}
App (Router + AuthContext)
├── ChildLayout          (vista niño, fondo neutro, sin header complejo)
│   ├── NavBar           (max 4 iconos grandes, siempre visible)
│   ├── PhraseBuilder    (Modulo A)
│   │   ├── PhraseBar
│   │   ├── CategoryStrip
│   │   ├── PictoGrid
│   │   │   └── PictoCard (x6 max) -- pictograma + WordLabel
│   │   ├── PlayButton    -- TTSEngine
│   │   └── ClearButton
│   ├── Anticipation     (Modulo B)
│   │   └── TimeSlotBlock (x3: Ahora / Despues / Luego)
│   └── EmotionRegulator (Modulo C)
│       ├── EmotionGrid   -- PictoCard x6
│       ├── IntensityPicker (opcional)
│       └── StrategyPanel -- PictoCard x3 grande
│
└── AdultLayout          (vista adulto, requiere PIN)
    ├── ChildSelector
    ├── DateRangePicker
    ├── KpiPanel
    │   └── KpiCard (x4)
    └── ChartsSection
        ├── LmfLineChart
        ├── CategoryBarChart
        └── EmotionDonutChart

-- Atoms compartidos entre layouts:
PictoImage  |  WordLabel  |  TTSEngine  |  PinGate
\end{lstlisting}
\end{cajatecnica}

% ══════════════════════════════════════════════════════════════════════════
%  8. HOJA DE RUTA
% ══════════════════════════════════════════════════════════════════════════
\section{Hoja de ruta de desarrollo}

\begin{tabular}{C{1.2cm} L{3.5cm} L{9cm}}
\toprule
\textbf{Fase} & \textbf{Periodo} & \textbf{Entregables} \\
\midrule
0 & Semanas 1--2 & Configuración de entorno, migraciones SQL, proxy ARASAAC \\
1 & Semanas 3--6 & Módulo A (PhraseBuilder) completo con TTS y registro \\
2 & Semanas 7--9 & Módulo C (EmotionRegulator) con logs \\
3 & Semanas 10--12 & Módulo B (Anticipation) con configuración adulto \\
4 & Semanas 13--16 & Panel adulto (KPIs + gráficas) \\
5 & Semanas 17--20 & Testing con familias piloto, ajustes UX, PWA offline \\
6 & Semana 21+ & Lanzamiento beta, revisión con profesionales sanitarios \\
\bottomrule
\end{tabular}

\begin{nota}[Criterio de priorización de módulos]
El Módulo A (frases) se implementa primero por ser el de mayor complejidad de estado y el que genera el mayor volumen de datos para los KPIs. El panel adulto se deja para cuando existan datos reales que visualizar.
\end{nota}

% ══════════════════════════════════════════════════════════════════════════
%  9. CONSIDERACIONES ÉTICAS Y DE PRIVACIDAD
% ══════════════════════════════════════════════════════════════════════════
\section{Consideraciones éticas y de privacidad}

Al tratarse de una aplicación que procesa datos de menores con discapacidad, se aplicarán las siguientes medidas desde el inicio del desarrollo:

\begin{itemize}
  \item \textbf{RGPD}: los datos de menores requieren consentimiento explícito del tutor legal. El backend no almacenará datos identificativos más allá del nombre y la fecha de creación del perfil.
  \item \textbf{Minimización de datos}: únicamente se registran los eventos necesarios para calcular los KPIs clínicamente relevantes.
  \item \textbf{Acceso por roles}: el PIN del adulto protege el panel de seguimiento. El niño no accede a datos de otros usuarios.
  \item \textbf{Sin publicidad}: la aplicación no mostrará publicidad ni compartirá datos con terceros.
  \item \textbf{Almacenamiento local primero}: los datos de sesión se guardarán localmente (IndexedDB) para uso offline, sincronizándose con el servidor cuando haya conexión.
  \item \textbf{Revisión por profesionales}: el sistema de KPIs está diseñado para ser interpretado junto a un logopeda o psicólogo, no como diagnóstico autónomo.
\end{itemize}

% ══════════════════════════════════════════════════════════════════════════
%  10. GLOSARIO
% ══════════════════════════════════════════════════════════════════════════
\section{Glosario}

\begin{tabular}{L{3.5cm} L{10.5cm}}
\toprule
\textbf{Término} & \textbf{Definición} \\
\midrule
CAA & Comunicación Aumentativa y Alternativa: conjunto de métodos y herramientas que complementan el habla. \\
TEA & Trastorno del Espectro Autista: condición del neurodesarrollo con afectación en comunicación social y comportamientos restringidos. \\
ARASAAC & Portal aragonés de CAA. Ofrece más de 40\,000 pictogramas bajo licencia Creative Commons. \\
LMF & Longitud Media de Frase: número promedio de elementos (pictogramas) por enunciado. \\
KPI & Key Performance Indicator: indicador clave de rendimiento o progreso. \\
PWA & Progressive Web App: aplicación web instalable en dispositivos sin necesidad de tiendas de apps. \\
TTS & Text-to-Speech: síntesis de voz a partir de texto. En esta app se usa la Web Speech API nativa. \\
JSONB & Tipo de dato de PostgreSQL que almacena JSON de forma binaria optimizada para consultas. \\
MVP & Minimum Viable Product: versión mínima funcional del producto para validación con usuarios reales. \\
PIN & Número de identificación personal usado para separar el acceso del niño y del adulto. \\
\bottomrule
\end{tabular}

% ══════════════════════════════════════════════════════════════════════════
%  FINAL
% ══════════════════════════════════════════════════════════════════════════
\vfill
\begin{center}
\textcolor{grisMedio}{\rule{0.6\linewidth}{0.4pt}}\\[6pt]
{\small\textcolor{gris}{Blueprint técnico preparado para fase de implementación.}}\\
{\small\textcolor{gris}{Versión 1.0 · 2025 · Documento interno}}
\end{center}

\end{document}